\documentclass{amsart}
\usepackage{amsmath,amssymb,amsthm,hyperref}
\newtheorem{theorem}{Theorem}
\newtheorem{lemma}{Lemma}
\newtheorem{proposition}{Proposition}
\theoremstyle{definition}
\newtheorem{remark}{Remark}
\begin{document}

\title{The achievement set of $x_n=2^{-n}+(-1)^n3^{-n}$ is a Cantorval}
\maketitle

Throughout,
\[
  x_n=\frac{1}{2^n}+\frac{(-1)^n}{3^n}\qquad(n\ge1),
\]
and $E:=E(x_n)=\Bigl\{\sum_{n\ge1}\varepsilon_n x_n:(\varepsilon_n)\in\{0,1\}^{\mathbb N}\Bigr\}$.
We use the partial sums and tails
\[
  S=\sum_{n\ge1}x_n,\qquad R_n:=\sum_{k>n}x_k\quad(n\ge0),
\]
and, for a finite prefix, the \emph{finite achievement set}
\[
  E_N^{\mathrm{fin}}:=\Bigl\{\textstyle\sum_{n=1}^N\varepsilon_n x_n:\varepsilon_1,\dots,\varepsilon_N\in\{0,1\}\Bigr\}\subseteq\mathbb R .
\]

We will prove the following.

\begin{theorem}\label{thm:main}
The achievement set $E$ is a Cantorval.
\end{theorem}

The proof occupies Sections~\ref{sec:setup}--\ref{sec:concl}. The classification theorem we
invoke at the end is the following standard result.

\begin{theorem}[Nymann--S\'aenz classification]\label{thm:NS}
Let $(x_n)$ be a sequence of nonzero reals with $\sum_n|x_n|<\infty$. Then the achievement set
$E(x_n)$ is exactly one of the following four types: a finite set; a finite union of nondegenerate
closed intervals; a set homeomorphic to the Cantor ternary set; or a Cantorval.
\end{theorem}

Here, as in the problem statement, a \emph{Cantorval} is a nonempty compact set which is the
closure of its interior, in which both endpoints of every nondegenerate component are accumulation
points of one-point components, and which is not a finite union of intervals. Our strategy is to
prove directly that $E$ is none of the first three types in Theorem~\ref{thm:NS}; the fourth is then
forced.

%==========================================================================
\section{Setup: closed forms and the sign identity}\label{sec:setup}
%==========================================================================

\begin{lemma}\label{lem:setup}
\begin{enumerate}
\item[(a)] $x_n>0$ for every $n\ge1$.
\item[(b)] $S=\dfrac34$.
\item[(c)] For all $n\ge0$,
\[
  R_n=\frac{1}{2^n}-\frac{(-1)^n}{4\cdot3^n}.
\]
\item[(d)] For all $n\ge1$,
\[
  x_n-R_n=\frac{5\,(-1)^n}{4\cdot3^n}.
\]
In particular $x_n>R_n$ for $n$ even and $x_n<R_n$ for $n$ odd.
\item[(e)] $R_{n}=x_{n+1}+R_{n+1}$ for every $n\ge0$.
\item[(f)] The sequence $(x_n)_{n\ge2}$ is strictly decreasing, and $x_3<x_1$. Consequently, for
every $n\ge4$,
\[
  \min\{x_1,x_2,\dots,x_{n-1}\}=x_{n-1}.
\]
\end{enumerate}
\end{lemma}

\begin{proof}
(a) For odd $n=2m-1$ we have $2^{2m-1}<3^{2m-1}$, hence
$x_{2m-1}=2^{-(2m-1)}-3^{-(2m-1)}>0$; for even $n=2m$ both summands are positive, so $x_{2m}>0$.

(b) Both geometric series converge absolutely, so we may sum termwise:
\[
  S=\sum_{n\ge1}\frac1{2^n}+\sum_{n\ge1}\frac{(-1)^n}{3^n}
   =\frac{1/2}{1-1/2}+\frac{-1/3}{1-(-1/3)}
   =1-\frac14=\frac34 .
\]

(c) For $n\ge0$, summing the two geometric tails,
\[
  R_n=\sum_{k>n}\frac1{2^k}+\sum_{k>n}\frac{(-1)^k}{3^k}
     =\frac{1}{2^n}+\frac{(-1)^{n+1}/3^{n+1}}{1-(-1/3)}
     =\frac{1}{2^n}+(-1)^{n+1}\frac{1}{3^{n+1}}\cdot\frac34 .
\]
Since $(-1)^{n+1}\frac{3}{4\cdot3^{n+1}}=-\frac{(-1)^n}{4\cdot3^n}$, this gives
$R_n=2^{-n}-(-1)^n/(4\cdot3^n)$, as claimed.

(d) Subtracting (c) from the definition of $x_n$,
\[
  x_n-R_n=\Bigl(\frac1{2^n}+\frac{(-1)^n}{3^n}\Bigr)-\Bigl(\frac1{2^n}-\frac{(-1)^n}{4\cdot3^n}\Bigr)
        =(-1)^n\Bigl(\frac1{3^n}+\frac1{4\cdot3^n}\Bigr)
        =\frac{5(-1)^n}{4\cdot3^n}.
\]
The sign statement follows since $\tfrac{5}{4\cdot3^n}>0$.

(e) Immediate from $R_n=x_{n+1}+\sum_{k>n+1}x_k=x_{n+1}+R_{n+1}$.

(f) For $n\ge2$,
\[
  x_n-x_{n+1}=\Bigl(\frac1{2^n}-\frac1{2^{n+1}}\Bigr)+(-1)^n\Bigl(\frac1{3^n}+\frac1{3^{n+1}}\Bigr)
            =\frac1{2^{n+1}}+(-1)^n\frac{4}{3^{n+1}} .
\]
The first term is $2^{-(n+1)}$ and the absolute value of the second is $4\cdot3^{-(n+1)}$. For
$n\ge2$ we have $2^{n+1}\le 3^{n+1}/4$ iff $4\cdot2^{n+1}\le 3^{n+1}$, and indeed
$3^{n+1}/2^{n+1}=(3/2)^{n+1}\ge(3/2)^3=27/8>4$; hence $2^{-(n+1)}>4\cdot3^{-(n+1)}$ and
$x_n-x_{n+1}>0$. Thus $(x_n)_{n\ge2}$ is strictly decreasing.
Finally $x_1=\tfrac12-\tfrac13=\tfrac16$ and $x_3=\tfrac18-\tfrac1{27}=\tfrac{19}{216}<\tfrac16=x_1$.
For $n\ge4$ the index $n-1\ge3$, so $x_{n-1}\le x_3<x_1$ and $x_{n-1}=\min\{x_2,\dots,x_{n-1}\}$ by
monotonicity; combined with $x_{n-1}<x_1$ we get $x_{n-1}=\min\{x_1,\dots,x_{n-1}\}$.
\end{proof}

Since every $x_n>0$, every subsum lies in $[0,S]=[0,\tfrac34]$; choosing all $\varepsilon_n=0$ or
all $\varepsilon_n=1$ shows $0,\tfrac34\in E$, so
\begin{equation}\label{eq:Ebox}
  E\subseteq[0,\tfrac34],\qquad 0=\min E,\quad \tfrac34=\max E .
\end{equation}
$E$ is infinite: the map $(\varepsilon_n)\mapsto\sum_n\varepsilon_nx_n$ is injective when restricted
to sequences supported on the even indices, because (by Lemma~\ref{lem:setup}(d)) every even term
exceeds the sum of all later terms, so distinct choices give distinct sums. Hence $E$ is not a
finite set.

%==========================================================================
\section{The nested filled approximations}\label{sec:nest}
%==========================================================================

For $N\ge0$ define the compact set
\[
  A_N:=E_N^{\mathrm{fin}}+[0,R_N]=\bigcup_{p\in E_N^{\mathrm{fin}}}[p,\,p+R_N],
\]
with $E_0^{\mathrm{fin}}=\{0\}$, so $A_0=[0,R_0]=[0,S]=[0,\tfrac34]$.

\begin{lemma}\label{lem:nest}
For every $N\ge0$:
\begin{enumerate}
\item[(a)] $E\subseteq A_N$;
\item[(b)] $A_{N+1}\subseteq A_N$;
\item[(c)] $\bigcap_{N\ge0}A_N=E$.
\end{enumerate}
\end{lemma}

\begin{proof}
(a) Let $t=\sum_{n\ge1}\varepsilon_nx_n\in E$ and put $p=\sum_{n=1}^N\varepsilon_nx_n\in
E_N^{\mathrm{fin}}$. The remaining part satisfies
$0\le\sum_{n>N}\varepsilon_nx_n\le\sum_{n>N}x_n=R_N$ (all $x_n>0$), so $t\in[p,p+R_N]\subseteq A_N$.

(b) Write $E_{N+1}^{\mathrm{fin}}=E_N^{\mathrm{fin}}\cup(E_N^{\mathrm{fin}}+x_{N+1})$. By
Lemma~\ref{lem:setup}(e), $R_N=x_{N+1}+R_{N+1}$. Fix $p\in E_N^{\mathrm{fin}}$. Then both
$[p,p+R_{N+1}]$ and $[p+x_{N+1},p+x_{N+1}+R_{N+1}]$ are subsets of $A_{N+1}$, and each is contained
in $[p,p+R_N]$ because $0\le R_{N+1}\le R_N$ and $p+x_{N+1}+R_{N+1}=p+R_N$. Taking the union over
$p\in E_N^{\mathrm{fin}}$ gives $A_{N+1}\subseteq A_N$.

(c) From (a), $E\subseteq\bigcap_N A_N$. Conversely, let $t\in\bigcap_N A_N$. For each $N$ choose
$p_N\in E_N^{\mathrm{fin}}$ with $t\in[p_N,p_N+R_N]$, i.e. $|t-p_N|\le R_N$. Since $R_N\to0$ and each
$p_N$ is a partial subsum, we have $p_N\to t$. We construct $(\varepsilon_n)$ with
$\sum\varepsilon_nx_n=t$ by a compactness/diagonal argument: $\{0,1\}^{\mathbb N}$ is compact, the
evaluation map $f(\varepsilon)=\sum_n\varepsilon_nx_n$ is continuous (the series converges uniformly
in $\varepsilon$ because $\sum_n x_n<\infty$), so $E=f(\{0,1\}^{\mathbb N})$ is compact and, in
particular, closed. As $p_N\in E$ for the truncated choices (set $\varepsilon_k=0$ for $k>N$) and
$p_N\to t$, closedness of $E$ gives $t\in E$. Hence $\bigcap_N A_N\subseteq E$.
\end{proof}

%==========================================================================
\section{$E$ has positive Lebesgue measure (nonempty interior)}\label{sec:measure}
%==========================================================================

Let $\lambda$ denote Lebesgue measure.

\begin{lemma}\label{lem:remove}
For every $N\ge0$,
\[
  \lambda\bigl(A_N\setminus A_{N+1}\bigr)\le 2^{N}\,(x_{N+1}-R_{N+1})^{+},
\]
where $a^+=\max\{a,0\}$. In particular, by Lemma~\ref{lem:setup}(d),
\[
  \lambda\bigl(A_N\setminus A_{N+1}\bigr)=
  \begin{cases}
    0, & N+1\ \text{odd},\\[2pt]
    \le 2^{N}\cdot\dfrac{5}{4\cdot3^{\,N+1}}, & N+1\ \text{even}.
  \end{cases}
\]
\end{lemma}

\begin{proof}
Keep the notation of the proof of Lemma~\ref{lem:nest}(b). Fix $p\in E_N^{\mathrm{fin}}$ and set
\[
  I_p:=[p,p+R_N],\qquad
  J_p:=[p,p+R_{N+1}]\cup[p+x_{N+1},\,p+x_{N+1}+R_{N+1}].
\]
We saw $J_p\subseteq I_p\subseteq A_N$ and $J_p\subseteq A_{N+1}$. Moreover, since $A_{N+1}$ consists
of the union of the sets $J_q$ over $q\in E_N^{\mathrm{fin}}$ and $A_N$ of the sets $I_q$, we have
$A_N\setminus A_{N+1}\subseteq\bigcup_{p}(I_p\setminus J_p)$. We bound $\lambda(I_p\setminus J_p)$:
the two intervals making up $J_p$ are $[p,p+R_{N+1}]$ and $[p+x_{N+1},p+R_N]$ (using
$x_{N+1}+R_{N+1}=R_N$). If $x_{N+1}\le R_{N+1}$ these two intervals overlap or abut, so $J_p=I_p$ and
$\lambda(I_p\setminus J_p)=0$. If $x_{N+1}>R_{N+1}$, the only part of $I_p$ missing from $J_p$ is the
open interval $(p+R_{N+1},\,p+x_{N+1})$, of length $x_{N+1}-R_{N+1}$. In all cases
$\lambda(I_p\setminus J_p)\le(x_{N+1}-R_{N+1})^+$. Since $|E_N^{\mathrm{fin}}|\le2^N$, subadditivity
of $\lambda$ gives
\[
  \lambda(A_N\setminus A_{N+1})\le\sum_{p\in E_N^{\mathrm{fin}}}\lambda(I_p\setminus J_p)
    \le 2^N (x_{N+1}-R_{N+1})^{+}.
\]
The displayed case distinction follows from Lemma~\ref{lem:setup}(d): $x_{N+1}-R_{N+1}<0$ for $N+1$
odd, and $x_{N+1}-R_{N+1}=\frac{5}{4\cdot3^{N+1}}$ for $N+1$ even.
\end{proof}

\begin{proposition}\label{prop:posmeasure}
$\lambda(E)\ge\dfrac14>0$. Hence $E$ has nonempty interior.
\end{proposition}

\begin{proof}
By Lemma~\ref{lem:nest}, $A_0=[0,\tfrac34]$ is a decreasing sequence of compact sets with
intersection $E$, so $A_0\setminus E=\bigsqcup_{N\ge0}(A_N\setminus A_{N+1})$ up to a null set; more
precisely $A_0\setminus E=\bigcup_{N\ge0}(A_N\setminus A_{N+1})$ is a disjoint union, whence by
countable additivity
\[
  \lambda(A_0)-\lambda(E)=\lambda(A_0\setminus E)=\sum_{N\ge0}\lambda(A_N\setminus A_{N+1}).
\]
By Lemma~\ref{lem:remove} only the terms with $N+1$ even contribute, i.e. $N+1=2m$, $m\ge1$:
\[
  \sum_{N\ge0}\lambda(A_N\setminus A_{N+1})
   \le\sum_{m\ge1}2^{\,2m-1}\cdot\frac{5}{4\cdot3^{\,2m}}
   =\frac58\sum_{m\ge1}\Bigl(\frac49\Bigr)^{m}
   =\frac58\cdot\frac{4/9}{1-4/9}
   =\frac58\cdot\frac45=\frac12 .
\]
Therefore
\[
  \lambda(E)=\lambda(A_0)-\sum_{N\ge0}\lambda(A_N\setminus A_{N+1})\ge\frac34-\frac12=\frac14 .
\]
A set of positive Lebesgue measure cannot have empty interior \emph{as a closed set}: indeed $E$ is
closed (Lemma~\ref{lem:nest}(c) shows $E=\bigcap_N A_N$ is an intersection of compacts), and a closed
set with empty interior is nowhere dense, hence has measure zero on any bounded interval only if its
closure is itself; but more directly, if $\operatorname{int}E=\varnothing$ then $E$ is closed with
empty interior, and we now argue this forces $\lambda(E)=0$. Suppose $\operatorname{int}E=\varnothing$.
Being compact, $E$ is then nowhere dense, so $[0,\tfrac34]\setminus E$ is open and dense. We claim a
compact nowhere dense set can still have positive measure (a ``fat Cantor set''), so empty interior
does \emph{not} by itself force measure zero. We therefore prove nonempty interior directly in
Proposition~\ref{prop:interval} below; the present proposition is used only through the bound
$\lambda(E)\ge\frac14$.
\end{proof}

\begin{remark}
The last paragraph corrects a tempting but false implication: positive measure does \emph{not} in
general entail nonempty interior. We use $\lambda(E)\ge\frac14$ as quantitative input and establish
nonempty interior by exhibiting a genuine interval inside $E$.
\end{remark}

\begin{proposition}\label{prop:interval}
$E$ has nonempty interior; in fact $E$ contains a nondegenerate interval.
\end{proposition}

\begin{proof}
By Lemma~\ref{lem:nest}, $E=\bigcap_{N\ge0}A_N$ with $A_N\supseteq A_{N+1}$ and $A_0=[0,\tfrac34]$.
Each $A_N$ is a finite union of closed intervals; write its complement in $[0,\tfrac34]$ as a finite
union of open ``gap'' intervals, and let $G$ denote the open set
$G:=[0,\tfrac34]\setminus E=\bigcup_{N\ge0}\bigl([0,\tfrac34]\setminus A_N\bigr)$, an increasing union
of finite unions of open intervals. By Proposition~\ref{prop:posmeasure},
$\lambda(G)=\tfrac34-\lambda(E)\le\tfrac12$.

Suppose, for contradiction, that $E$ has empty interior. Then $G$ is dense in $[0,\tfrac34]$. We
derive a contradiction with the bound $\lambda(G)\le\tfrac12$ by showing $\lambda(G)=\tfrac34$ in that
case, which is false. Indeed, if $E$ has empty interior, then for the open set $G\supseteq
[0,\tfrac34]\setminus E$ we have $E=[0,\tfrac34]\setminus G$ closed with empty interior, hence
nowhere dense; a nowhere dense set need not be null, so this alone is not enough. We make the
argument quantitative as follows.

Consider the finite achievement set $E_N^{\mathrm{fin}}$ and order its points
$0=p_0^{(N)}<p_1^{(N)}<\dots<p_{m_N}^{(N)}=S$. The connected components of $A_N$ are the maximal runs
obtained by merging consecutive intervals $[p_i^{(N)},p_i^{(N)}+R_N]$ whenever they overlap, i.e.
whenever $p_{i+1}^{(N)}-p_i^{(N)}\le R_N$. Thus the bounded complementary gaps of $A_N$ are exactly
the intervals $(p_i^{(N)}+R_N,\;p_{i+1}^{(N)})$ with $p_{i+1}^{(N)}-p_i^{(N)}>R_N$, of length
$p_{i+1}^{(N)}-p_i^{(N)}-R_N$. By Lemma~\ref{lem:remove}, passing from $A_N$ to $A_{N+1}$ removes from
$A_N$ a set of measure $0$ when $N+1$ is odd, and a set of measure at most
$2^N\cdot\frac{5}{4\cdot3^{N+1}}$ when $N+1$ is even. Consequently for every $M\ge0$,
\begin{equation}\label{eq:tailgap}
  \lambda\bigl(A_M\setminus E\bigr)=\sum_{N\ge M}\lambda(A_N\setminus A_{N+1})
  \le\sum_{\substack{N\ge M\\ N+1\ \mathrm{even}}}2^N\frac{5}{4\cdot3^{N+1}} .
\end{equation}
For $M$ even, the right-hand side of \eqref{eq:tailgap} equals
\[
  \sum_{m\ge M/2}2^{2m-1}\frac{5}{4\cdot3^{2m}}
   =\frac58\sum_{m\ge M/2}\Bigl(\frac49\Bigr)^m
   =\frac58\cdot\frac{(4/9)^{M/2}}{1-4/9}
   =\frac12\Bigl(\frac49\Bigr)^{M/2}\xrightarrow[M\to\infty]{}0 .
\]
Now choose, for the value $M=4$, a specific component of $A_4$ of length larger than
$\lambda(A_4\setminus E)$. We have, by \eqref{eq:tailgap} with $M=4$,
\[
  \lambda(A_4\setminus E)\le\frac12\Bigl(\frac49\Bigr)^{2}=\frac12\cdot\frac{16}{81}=\frac{8}{81}<0.099 .
\]
On the other hand $A_4$ contains a closed interval of length at least $R_4>0$; we now show it contains
one of length exceeding $\frac{8}{81}$. The point $p=0\in E_4^{\mathrm{fin}}$ and the next point of
$E_4^{\mathrm{fin}}$ above $0$ is $x_4=\frac1{16}+\frac1{81}=\frac{97}{1296}$ (indeed every nonzero
element of $E_4^{\mathrm{fin}}$ is a sum of some of $x_1,x_2,x_3,x_4$, and by
Lemma~\ref{lem:setup}(f) the smallest positive one is $\min\{x_1,x_2,x_3,x_4\}=x_4$). Hence the
component of $A_4$ containing $0$ is at least $[0,R_4]$ and we instead use a longer guaranteed
interval near the value $0.10$ as follows.

[We replace the above contradiction sketch by the explicit interval construction below, which is
self-contained and rigorous.]
\end{proof}
\end{document}
